\documentclass[a4paper]{article}
% Español (México) con XeLaTeX: fontspec para fuentes Unicode, polyglossia
% para la división silábica y los nombres del documento en la variante mexicana.
\usepackage{fontspec}
\usepackage{polyglossia}
\setdefaultlanguage[variant=mexican]{spanish}
% Los nombres chinos van como «pinyin (original)»: xeCJK da a los caracteres
% Han su propia fuente; sin ella XeLaTeX los omite con un aviso.
\usepackage{xeCJK}
\setCJKmainfont{FandolSong-Regular.otf}
% polyglossia no traduce los nombres de listings.
\AtBeginDocument{\providecommand{\lstlistingname}{}\renewcommand{\lstlistingname}{Listado}%
  \providecommand{\lstlistlistingname}{}\renewcommand{\lstlistlistingname}{Índice de listados}}
\usepackage{amsmath, amssymb}
\usepackage{graphicx}
\usepackage[margin=2.5cm]{geometry}
\usepackage[most]{tcolorbox}
\usepackage{etoolbox}
\usepackage{listings}
\usepackage{hyperref}
\usepackage{booktabs}
\usepackage{subcaption}
\usepackage{float}
\newtcolorbox{knowledgebox}[1]{enhanced, colback=blue!5!white, colframe=blue!75!black, colbacktitle=blue!75!black, coltitle=white, fonttitle=\bfseries, title={#1}, attach boxed title to top left={yshift=-2mm, xshift=2mm}, boxrule=1pt, sharp corners}
\newtcolorbox{importantbox}[1]{enhanced, colback=yellow!10!white, colframe=yellow!80!black, colbacktitle=yellow!80!black, coltitle=black, fonttitle=\bfseries, title={#1}, sharp corners}
\newtcolorbox{warningbox}[1]{enhanced, colback=red!5!white, colframe=red!75!black, colbacktitle=red!75!black, coltitle=white, fonttitle=\bfseries, title={#1}, sharp corners}
\lstset{language=Python, basicstyle=\ttfamily\small, keywordstyle=\color{blue}, stringstyle=\color{red!60!black}, commentstyle=\color{green!60!black}, breaklines=true, frame=single, numbers=left, numberstyle=\tiny\color{gray}, captionpos=b, extendedchars=true}

\newcommand{\notetitle}{Perspectiva geométrica de RoPE y detalles del cálculo de RoPE en Qwen3}
\newcommand{\noteauthors}{Elaboradas a partir de materiales públicos del curso}
\newcommand{\notedate}{2025}
\newcommand{\videochannel}{Wudaokou Nash (五道口纳什)}
\newcommand{\videopublishdate}{2025}
\newcommand{\videoduration}{aproximadamente 21 minutos}
\newcommand{\videourl}{}
\newcommand{\videocoverpath}{cover.jpg}
\newcommand{\repourl}{https://github.com/hqhq1025/ai-course-notes}


% Footer with repo info
\usepackage{fancyhdr}
\pagestyle{fancy}
\fancyhf{}
\fancyfoot[C]{\thepage}
\fancyfoot[R]{\footnotesize\color{gray}\href{\repourl}{AI Course Notes}}
\renewcommand{\headrulewidth}{0pt}
\renewcommand{\footrulewidth}{0pt}

\begin{document}
\begin{titlepage}
\centering
{\Large Notas del curso\par}\vspace{1.2cm}
{\huge\bfseries \notetitle\par}\vspace{0.8cm}
{\large \noteauthors\par}\vspace{0.3cm}
{\large \notedate\par}\vspace{1.2cm}
\ifdefempty{\videocoverpath}{}{\includegraphics[width=0.82\textwidth,height=0.45\textheight,keepaspectratio]{\videocoverpath}\par}
\vfill
\begin{tcolorbox}[width=0.9\textwidth, colback=black!2!white, colframe=black!60, sharp corners]
\textbf{Autor o canal del video}: \videochannel\par
\textbf{Fecha de publicación}: \videopublishdate\par
\textbf{Duración del video}: \videoduration\par
\ifdefempty{\videourl}{}{\textbf{Enlace del video}: \href{\videourl}{\nolinkurl{\videourl}}\par}\par
\textbf{Página del proyecto}: \href{\repourl}{\nolinkurl{\repourl}}\par
\end{tcolorbox}
\end{titlepage}
\tableofcontents
\newpage

\section{Introducción: la insensibilidad a la posición del Multi-Head Self-Attention}

El Multi-Head Self-Attention es en sí mismo \textbf{insensible a la posición}. Si no se introduce una codificación de posición, equivale a un modelo de bolsa de palabras (bag of words) más capaz. RoPE (Rotary Position Embedding) inyecta con ingenio la información de posición en el cálculo de la atención mediante una matriz de rotación.

\begin{importantbox}{Idea central de RoPE}
Al asignar a $Q$ y a $K$ ángulos de rotación relacionados con su \textbf{posición absoluta}, se aprovechan las propiedades matemáticas de la matriz de rotación para lograr, en última instancia, una \textbf{codificación de posición relativa}. Esta relatividad no se construye de antemano, sino que se deriva de manera natural en las matemáticas.
\end{importantbox}

\section{Expresión matemática de RoPE}

\subsection{Propiedades básicas de la matriz de rotación}

Se aplica la matriz de rotación $R_m$ al Query de la posición $m$, y $R_n$ al Key de la posición $n$:

$$
\text{score}(m, n) = (R_m q)^T (R_n k) = q^T R_m^T R_n k = q^T R_{n-m} k
$$

Propiedad clave: $R_m^T R_n = R_{n-m}$, la transpuesta de la matriz de rotación es igual a la rotación inversa, y la combinación de dos rotaciones solo depende de la diferencia de posiciones $n - m$.

\begin{warningbox}{La relatividad se deduce, no se construye}
Solo necesitamos calcular de antemano la matriz de rotación para cada posición absoluta; la codificación de posición relativa es un resultado natural de la deducción matemática y no requiere construirse de forma explícita.
\end{warningbox}

\subsection{Parámetro de frecuencia $\theta$}

RoPE aplica rotaciones de distinta frecuencia sobre los diferentes subespacios 2D (subspace) del vector de embedding:

$$
\theta_i = \text{base}^{-2i/d}
$$

\begin{itemize}
  \item $\text{base}$: generalmente se toma 10,000 (RoPE estándar); Qwen3 toma \textbf{1,000,000} (para lograr una ventana de contexto mayor)
  \item $i$: índice del subespacio, de 0 a $d/2 - 1$
  \item $d$: head dimension
\end{itemize}

$\theta_i$ decae de forma \textbf{exponencial} conforme aumenta $i$; al tomar $\log$, decae linealmente.

\subsection{Resumen de la sección}

RoPE logra la codificación de posición relativa mediante matrices de rotación de posiciones absolutas; $\theta_i$ controla la frecuencia de rotación de cada subespacio. Cuanto mayor sea base, más lenta será la rotación de los subespacios de baja frecuencia, lo cual favorece el modelado de largo alcance.
\section{Perspectiva geométrica: entender RoPE desde la frecuencia}

\subsection{Propiedad multiescala}

El decaimiento exponencial de $\theta_i$ implica que RoPE tiene una propiedad \textbf{multiescala}:

\begin{importantbox}{El espectro de RoPE}
\begin{itemize}
  \item \textbf{Índice bajo ($i$ pequeño)}: $\theta_i$ es grande, la frecuencia de rotación es alta $\to$ codifica dependencias de \textbf{corto alcance} (sintaxis, combinaciones de palabras como ``New York'')
  \item \textbf{Índice alto ($i$ grande)}: $\theta_i$ es pequeño, la frecuencia de rotación es baja $\to$ codifica dependencias de \textbf{largo alcance} (semántica, estructura de párrafo, referencias remotas)
\end{itemize}
\end{importantbox}

\subsection{La configuración de RoPE en Qwen3}

El RoPE estándar usa base = 10000, mientras que la serie Qwen3 usa base = 1000000. Una base mayor implica que todas las frecuencias son más bajas, la rotación del subespacio de baja frecuencia es más lenta, lo cual favorece capturar dependencias de mayor alcance y sustenta una context window más grande.

\subsection{Comparación con el esquema interleaved estándar}

Al presentar antes LLaMA se usó la versión interleaved ingenua de RoPE (que trata las dimensiones adyacentes como un subespacio 2D). Es posible que la implementación de Qwen3 difiera en la forma de dividir el subespacio; hace falta revisar el código concreto para confirmarlo.

\subsection{Resumen de la sección}

La esencia geométrica de RoPE consiste en aplicar rotaciones a distintas velocidades sobre los diferentes subespacios 2D del embedding. El subespacio de alta frecuencia se encarga de las relaciones sintácticas de corto alcance, y el subespacio de baja frecuencia se encarga de las relaciones semánticas de largo alcance. Qwen3 sustenta una ventana de contexto más larga aumentando la base.
\section{Indicaciones de implementación: el contexto largo no depende solo de un base}
\subsection{Cómo afecta la configuración de RoPE al desempeño de ingeniería}
El video dejó al final una indicación de ingeniería muy importante: cambiar el base de 10,000 a 1,000,000 es solo una parte de la estrategia de contexto largo. Que el modelo pueda realmente aprovechar de forma estable un contexto más largo también depende de la distribución de longitudes vista durante el entrenamiento, de la implementación de la caché KV y de la estabilidad numérica de la atención a distancias muy largas.

\begin{table}[H]
\centering
\begin{tabular}{p{3.2cm}p{4.6cm}p{4.6cm}}
\toprule
Factor & Efecto directo & Implicación de ingeniería \\
\midrule
Base de RoPE & Cambia la frecuencia de rotación de cada subespacio & Determina la resolución de las diferencias de posición a larga distancia \\
Distribución de longitudes de entrenamiento & Determina si el modelo realmente ha visto contexto largo & Las longitudes no vistas suelen seguir degradándose \\
Caché KV / implementación de inferencia & Determina el costo y la estabilidad de la inferencia con secuencias largas & Afecta si la ventana teórica realmente se puede ejecutar \\
\bottomrule
\end{tabular}
\caption{Además de los parámetros de RoPE, la capacidad de contexto largo también está determinada por varios factores de ingeniería en conjunto}
\end{table}

\begin{knowledgebox}{Por qué vale la pena entender la geometría}
Una vez que se entiende RoPE desde la perspectiva geométrica y espectral, es más fácil determinar si un problema de contexto largo se debe a que «la codificación misma no es suficientemente buena» o a que «el entrenamiento y la inferencia no liberaron realmente el potencial de la codificación». Esta es también la razón por la que muchos reportes de modelos discuten a la vez RoPE, la distribución de longitud de los datos y la optimización de la inferencia.
\end{knowledgebox}

\subsection{Resumen de la sección}
La fórmula de RoPE explica cómo se codifica la relación de posición dentro de la atención, pero la capacidad real de contexto largo sigue siendo el producto conjunto de «codificación + entrenamiento + implementación de inferencia». Entender esto ayuda a interpretar correctamente las decisiones de configuración de modelos como Qwen3.
\section{Síntesis y ampliación}

\begin{enumerate}
  \item RoPE implementa la codificación de posición relativa mediante rotaciones de posición absoluta; la relatividad es resultado de la deducción matemática
  \item $\theta_i = \text{base}^{-2i/d}$ describe un espectro multiescala que va de alta a baja frecuencia
  \item Subespacio de alta frecuencia $\to$ dependencias de corto alcance; subespacio de baja frecuencia $\to$ dependencias de largo alcance
  \item Qwen3 usa base = 1000000, lo que admite una ventana de contexto más grande
  \item En la siguiente sesión se deducirá la fórmula de decaimiento a distancia del attention score
\end{enumerate}

\subsection{Lecturas adicionales}
\begin{itemize}
  \item Su et al., ``RoFormer: Enhanced Transformer with Rotary Position Embedding'' (2021)
  \item Explicación sobre la configuración de RoPE en el reporte técnico de Qwen3
  \item Video previo de la serie: introducción básica a RoPE en LLaMA 2
\end{itemize}

\end{document}
